\documentclass[10pt,a4paper]{article}
\usepackage[margin=2cm]{geometry}
\usepackage{booktabs,tabularx,array,longtable}
\usepackage{amsmath}
\usepackage{xcolor}
\usepackage{enumitem}
\usepackage{hyperref}
\hypersetup{colorlinks=true,linkcolor=blue!50!black,urlcolor=blue!50!black}
\setlist{nosep,leftmargin=1.4em}
\newcolumntype{L}{>{\raggedright\arraybackslash}X}

\title{RadarNet-1D: Changes to Software and Hardware\\ and Response to the VLSID 2027 Reviews (Submission \#119)}
\author{Prepared for the authors}
\date{26 October 2026}

\begin{document}
\maketitle

\section*{Summary}
The original accelerator was re-examined before revising the paper. Its RTL did not compute
the trained network correctly, and its reported latency was misread from the
waveform. The hardware was therefore redesigned (\texttt{rtl\_v2/}), an integer-only
quantization scheme with a bit-exact golden model was introduced, and the design was
verified at four levels:
\begin{enumerate}
  \item in Python;
  \item in RTL simulation (Icarus Verilog and Vivado xsim);
  \item in Vivado implementation (setup and hold timing closed);
  \item on the ZedBoard.
\end{enumerate}
On the board it reproduces the golden model bit-exactly on 50/50 test vectors and is
$60.5\times$ faster than the same model on the on-chip ARM Cortex-A9. Every number in the
revised paper is taken from these runs. Raw Vivado reports are in
\texttt{results/vivado\_reports/}, and the board measurement is in
\texttt{results/zedboard\_measured.md}.

\begin{center}
\small
\begin{tabularx}{\textwidth}{l l L}
\toprule
\textbf{Quantity} & \textbf{Submitted paper} & \textbf{Verified now} \\
\midrule
Cycles per inference & 56 & 20,315 (simulation); 23,013 on board incl.\ AXI input \\
Latency & 0.78~$\mu$s & 142.2~$\mu$s compute; 161.9~$\mu$s measured end to end \\
Clock & 72.2~MHz (then 100) & 142.86~MHz system (met); 175.4~MHz out of context (met) \\
FP32 accuracy & 83.23\% & 83.26\% (\texttt{metrics.json}, reproduced in Colab) \\
INT8 accuracy & 83.18\% (weights only) & 83.18\% integer-only, executed bit-exactly by HW \\
HW correctness & not verified & 50/50 bit-exact (RTL), 50/50 bit-exact (board) \\
LUT / FF / DSP / BRAM & 10,012 / 5,461 / 0 / 15 & 3,117 / 2,448 / 104 / 14 (whole PL) \\
Power & 0.671~W (vectorless) & 0.290~W accelerator, 1.978~W chip incl.\ PS7 (vectorless) \\
Baseline & GPU (estimated) & ARM Cortex-A9, same chip, measured: 9.79~ms \\
\bottomrule
\end{tabularx}
\end{center}

\section{Problems Found in the Submitted Design}
\begin{enumerate}
  \item \textbf{Latency.} The original \texttt{conv1d\_engine} performs one MAC every two cycles,
  about 1.18~M MACs in total. Its own waveform (Fig.~3 of the submitted paper) shows
  \texttt{total\_cycles = 0x244452} = 2,376,786 cycles, i.e.\ 32.9~ms at 72.2~MHz. The
  value ``56 cycles'' was read from the cycle counter during the LOAD state.
  \item \textbf{No requantization.} Each 24-bit accumulator was saturated directly to INT8, with
  biases quantized at unrelated scales. The hardware therefore did not compute the trained
  network: in simulation it matched 0 of 4 expected classes, and the submitted Fig.~3 itself
  shows wrong classes (vector~0: 3 instead of 2; vector~1: 1 instead of 0).
  \item \textbf{Memory timing and addressing.} The synchronous block-RAM read latency was not
  accounted for, so each MAC used the previous weight. The RB2 projection layer read the
  weights of RB2C1 from the ROM.
  \item \textbf{Residual addition.} The two INT8 operands were added without aligning their scales.
  \item \textbf{Accuracy claim.} The 83.18\% INT8 figure came from rounding the weights only, with
  FP32 activations, which is not the arithmetic of the hardware.
  \item \textbf{Text.} The paper claimed DSP48E1 MAC engines (Table V: 0 DSPs), a single 36~KB BRAM
  (15 used), and an ``adaptive-precision'' accelerator.
\end{enumerate}

\section{Software Changes}
{\small
\begin{tabularx}{\textwidth}{>{\raggedright\arraybackslash\ttfamily}p{4.5cm} L}
\toprule
\textrm{\textbf{File}} & \textbf{Change} \\
\midrule
python/int\_quant.py & \textit{New.} Integer-only quantization and bit-exact golden model.
  It folds BN, quantizes weights per channel (INT8), calibrates per-tensor activation scales
  (99.99th percentile, 4,096 training frames), and computes INT32 biases, fixed-point
  requantization multipliers $M<2^{15}$ with per-layer shifts, and scale-aligned residual
  constants $M_s 2^{SK}$. Sub-commands \texttt{calib}, \texttt{eval}, \texttt{export}
  (ROM images and layer-descriptor table for any $W$), and \texttt{export-c} (C headers). \\
python/preprocess.py & Streaming loader for the 21~GB HDF5 file (about 2~GB of RAM instead of
  about 25~GB). Verified to produce identical arrays and therefore the identical train/test
  split as before. \\
sw/rn\_infer.c & \textit{New.} The same integer model in C (ARM baseline). Bit-exact with the
  golden model on 50/50 vectors (4-class and 7-class models). \\
sw/rn\_bench.c & \textit{New.} Bare-metal benchmark: runs every test vector on the ARM and on the
  accelerator over AXI, checks both against the golden logits, times both, and prints inputs,
  logits, predictions and a summary on the UART. \\
sw/build\_vitis.py & \textit{New.} Scripted Vitis platform and application build
  (Cortex-A9 / A53, \texttt{-O3}). \\
colab/*.py & \textit{New.} One-file Colab scripts. Step 1 downloads RadioML, rebuilds the split,
  calibrates and evaluates the integer model. Step 2 runs the Option A pipeline. \\
python/anomaly\_dataset.py\newline train\_anomaly.py\newline baselines.py\newline anom\_eval.py\newline anom\_report.py &
  \textit{New (Option A, not yet in the paper).} RadioML-based jamming / spoofing / interference
  dataset with controlled jammer-to-signal ratio (JSR), training, an energy-detector baseline,
  an expert-feature classifier, and ROC / $P_d$-at-fixed-$P_{fa}$ evaluation. \\
\bottomrule
\end{tabularx}}

\section{Hardware Changes}
\subsection{Architecture (\texttt{rtl\_v2/src/})}
\begin{tabularx}{\textwidth}{l L}
\toprule
\textbf{Original module} & \textbf{Replaced by} \\
\midrule
\texttt{conv1d\_engine} (one per layer, serial) & One shared \texttt{rn\_pe\_array}: $W{\times}W$ INT8 multipliers (DSP48) with registered adder trees, $W^2$ MACs per cycle. \\
\texttt{proj\_conv}, \texttt{fc\_layer} & Ordinary layers on the same engine (1$\times$1 stride-2 projection; FC with $L{=}1$). \\
Top-level 8-state FSM & Descriptor-driven sequencer in \texttt{rn\_top}: loop nest $o_g \rightarrow t \rightarrow i_g \rightarrow k$, layer table generated by Python. \\
Skip-add in FSM, \texttt{relu} & Six-stage requantization pipeline R0--R5: bias, $\times M$, $+$ shortcut$\cdot M_s 2^{SK}$, rounding, shift, ReLU, saturation. \\
8 per-layer output RAMs & Three ping-pong feature buffers (\texttt{rn\_buf}), word $=W$ channels. \\
\texttt{weight\_rom} (BRAM IP, 8 bit) & Wide weight ROM ($W{\times}W$ weights per row) and parameter ROM (INT32 bias, $M$ per channel). \\
\texttt{gap\_unit} & Lane-parallel \texttt{rn\_gap} with rounding shift. \\
\texttt{argmax\_out} (4 classes) & Sequential argmax over any number of classes (7-class model verified). \\
--- & \texttt{rn\_axi}: AXI4-Lite wrapper (register map, debug ports for ILA/VIO); ZedBoard and Zynq UltraScale+ system scripts. \\
\bottomrule
\end{tabularx}

\subsection{Timing closure (pipelining)}
At a 5~ns target the first version reached about 105~MHz. Each round below was re-verified
bit-exact before the next.
\begin{enumerate}
  \item Two-stage argmax; reset synchronizer.
  \item Registered GAP datapath.
  \item An R0 register stage before the requantization multipliers.
  \item Registered layer descriptors.
  \item DSP48 mapping of the PE multipliers, and a separate shift stage.
  \item Two-cycle BRAM reads (output registers) and two-stage address generation.
  \item Pre-scaled shortcut constant, so the DSP operand comes straight from a register;
  equality loop-bound compares.
\end{enumerate}
\textbf{Result.} Internal setup and hold slack are met at 169.5 / 175.4 / 138.9~MHz for
$W=4/8/16$. These frequencies are confirmed by re-implementing each design at that frequency,
with the clock on a global buffer. The full ZedBoard system meets all constraints at 142.86~MHz
(WNS $+0.385$~ns, WHS $+0.053$~ns, WPWS $+2.520$~ns).

\subsection{Verification}
\begin{itemize}
  \item \textbf{RTL simulation.} All four INT32 logits match the golden model on 50/50 vectors
  for $W$ = 4, 8, 16 and 32, for the 4-class and the 7-class model, in Icarus Verilog and in
  Vivado xsim.
  \item \textbf{AXI wrapper.} Verified with an AXI4-Lite testbench. An initial
  handshake-ordering bug was found and fixed.
  \item \textbf{Board.} ZedBoard, user-built bitstream and Vitis application: 50/50 bit-exact.
  Cortex-A9 9,787~$\mu$s vs.\ accelerator 161.9~$\mu$s end to end ($60.5\times$).
\end{itemize}

\section{Response to the Reviewers}
We thank the reviewers for their careful reading and constructive comments. Re-examining the
design in response to them showed that the latency reported in the submission was wrong and
that the original RTL did not execute the trained network correctly. The accelerator was
therefore redesigned and re-verified, and the revised manuscript reports only values obtained
from simulation, Vivado implementation, and measurement on the ZedBoard. Our responses follow
each comment, with references to the revised manuscript.

\newcommand{\rc}[2]{\par\medskip\noindent\textbf{#1}\par\noindent\textit{#2}\par\smallskip\noindent\textbf{Response:} }

\subsection*{Reviewer 1}

\rc{Comment 1.1 --- Incremental novelty}{BN folding, INT8 quantization, GAP through a shift, residual blocks, MAC engines, BRAM weight storage, and FSM-controlled RTL are all established techniques. The authors should clarify what hardware technique is genuinely new beyond combining these established components.}
We agree that these building blocks are individually well established, and we no longer present them as contributions. The contributions of the revised paper (Sec.~I) are:
\begin{enumerate}
  \item an integer-only INT8 scheme for a residual 1D CNN, with per-channel weight scales, 32-bit
  biases, fixed-point requantization, and \emph{scale-aligned} residual additions (Eq.~6). Its
  bit-exact golden model is used for verification from Python through RTL simulation to the
  physical board;
  \item a parameterizable $W{\times}W$ processing-element architecture in which a single PE array
  is time-multiplexed over all layers by a descriptor-driven sequencer, so the network is
  defined by a generated table rather than by per-layer hardware. A design-space exploration
  over $W$ with timing-closed frequencies quantifies the latency--resource trade-off
  (Table~VI);
  \item on-board validation: bit-exact results on silicon, and a measured comparison against the
  same model on the device's own ARM processor (Table~IX).
\end{enumerate}

\rc{Comment 1.2 --- Class detection}{How are the 24 RadioML modulation classes mapped to Normal, Jammer, Spoofer, and Interference? Also, please share how modulation type can reliably identify these operational categories, particularly Jammer.}
The complete mapping is now given in Table~I. We agree that a modulation type alone does not
establish that a signal is a jammer, and the revised text states this explicitly. The grouping
is presented as a candidate operational mapping used to demonstrate edge-deployed signal
categorization, and its validation against operational signal libraries is stated as future
work. In addition, we have built a RadioML-based dataset in which actual interference waveforms
are injected into RadioML signals at controlled jammer-to-signal ratios:
\begin{itemize}
  \item barrage, tone, swept, and pulsed jamming;
  \item replay spoofing;
  \item co-channel interference.
\end{itemize}
It is evaluated against energy-detector and expert-feature baselines. The accelerator runs this
seven-class model without hardware changes; it has been verified bit-exactly in RTL simulation.
This extension is outlined as future work in the Conclusion.

\rc{Comment 1.3 --- Inference latency}{The accelerator loads 128 I/Q pairs but claims only 56 cycles for complete inference. How are the 128 samples loaded within this latency?}
The reviewer is correct: the 56-cycle figure was wrong. It was read from the cycle counter during the
LOAD state, and the original serial datapath in fact required about 2.38 million cycles. The
redesigned accelerator performs 64 MACs per cycle. One inference takes 20,315 cycles from start
to result in RTL simulation, including the 128-cycle frame load, i.e.\ 142.2~$\mu$s at
142.86~MHz. On the ZedBoard, the samples are written by the ARM processor over AXI4-Lite and the
cycle counter reads 23,013 cycles; the measured end-to-end latency, including this transfer, is
161.9~$\mu$s (Sec.~IV-B, Table~IX).

\rc{Comment 1.4 --- Comparison with the state of the art}{The proposed accelerator should be quantitatively compared with the prior FPGA-based RF anomaly detector by Moss et al., which also targets real-time FPGA anomaly detection.}
The work of Moss et al.~[10] uses an unsupervised autoencoder that flags out-of-distribution
signals, whereas our design is a supervised classifier that assigns operational categories. The
two tasks and their metrics differ, so the comparison in the revised introduction (Sec.~I)
contrasts the two approaches. A quantitative entry for [10] will be added to the comparison
table (Table~V), reporting its published latency, throughput and resource figures alongside
ours.

\subsection*{Reviewer 2}

\rc{Comment 2.1 --- Algorithmic description}{There is no algorithmic description of the proposed methodology.}
A new subsection (Sec.~III-A) describes the complete integer inference algorithm:
\begin{itemize}
  \item BN folding (Eqs.~3--4);
  \item per-channel weight quantization (Eq.~5);
  \item activation calibration and INT32 bias computation;
  \item the requantization and residual-addition step executed by the hardware (Eq.~6), with the
  accumulator-overflow bound.
\end{itemize}
Sec.~III-B describes the hardware algorithm: the loop nest (output-channel group, time step,
input-channel group, kernel tap), the layer sequence, and the buffer allocation.

\rc{Comment 2.2 --- Block architecture}{There is no overall block/architecture of the proposed hardware.}
Two new figures describe the hardware:
\begin{itemize}
  \item Fig.~2 shows the top-level datapath: sequencer, three feature buffers, weight and parameter
  ROMs, PE array, accumulators, requantization pipeline with residual input, GAP unit, and
  argmax;
  \item Fig.~3 details one PE lane: the DSP48 multipliers, the registered adder tree, the
  accumulator, and the six-stage requantization pipeline.
\end{itemize}
The integration with the Zynq processing system (AXI4-Lite register interface, clocking,
reset) is described in Sec.~III-B.

\rc{Comment 2.3 --- Comparison with recent work}{Comparisons with the recent state-of-the-art publications are missing.}
The comparison table (Table~V) has been revised. The entries for our design now report verified
values: accuracy of the integer model, compute and measured latency, accelerator power,
resources, parameter storage, and on-board validation. The design-space exploration (Table~VI)
and the measured ARM comparison (Table~IX) add design-internal and same-platform comparisons.
The discussion now positions the design against streaming and binarized accelerators by
operating point, rather than by latency alone.

\subsection*{Reviewer 3}

\rc{Comment 3.1 --- Novelty}{The novelty of the proposed work appears to be limited, as the techniques employed are well-established and commonly used in similar FPGA-based accelerator architectures. The authors should clearly identify and substantiate the specific novel contributions.}
Please see our response to Comment~1.1. The contributions are restated in Sec.~I, and each is
substantiated by results: Table~IV (integer-only accuracy), Table~VI (design-space exploration
with timing closure), and Table~IX (bit-exact on-board measurement).

\rc{Comment 3.2 --- Operating frequency}{What is the rationale for selecting an operating frequency of 72.2~MHz? The authors should justify this frequency and explain whether it is determined by the FPGA implementation constraints or by a specific architectural consideration.}
The 72.2~MHz figure was the maximum frequency estimated from a failing 100~MHz implementation of the
original design, which had long combinational paths. The redesigned datapath is pipelined
throughout:
\begin{itemize}
  \item registered multipliers and adder-tree levels;
  \item BRAM output registers;
  \item two-stage address generation;
  \item a six-stage requantization pipeline;
  \item registered control.
\end{itemize}
Out of context, the $W{=}8$ accelerator closes timing (setup and hold) at 175.4~MHz. The
system clock of 142.86~MHz is the fastest Zynq PS clock step (1000/7~MHz) at which the
complete system, including the AXI interconnect, meets timing with margin: WNS $+0.385$~ns,
WHS $+0.053$~ns (Sec.~IV-A).

\rc{Comment 3.3 --- LOAD and inference cycles}{The paper reports that the LOAD operation and inference can be completed within 56 clock cycles. This appears to be very aggressive and requires further clarification.}
Please see our response to Comment~1.3. The corrected figures are 20,315 cycles in simulation
and 161.9~$\mu$s measured on the board, both including the input load.

\rc{Comment 3.4 --- FPGA-SoC block diagram}{A detailed FPGA-SoC block diagram is not provided. Including a detailed block diagram showing the processor, accelerator, memory hierarchy, interconnects, and other major components, along with relevant implementation details, would provide better clarity on the overall system architecture and implementation.}
A system block diagram has been prepared. It shows:
\begin{itemize}
  \item the ARM Cortex-A9 processing system, DDR and on-chip memory;
  \item the AXI interconnect and the AXI4-Lite accelerator interface with its register map;
  \item the accelerator, the clock and reset generation, and the JTAG host connection.
\end{itemize}
The accelerator's internal memory hierarchy (three feature buffers, weight and parameter ROMs)
is shown in Fig.~2. Implementation details of the system are given in Sec.~III-B and
Tables~VII--VIII.

\subsection*{Reviewer 4}

\rc{Comment 4.1 --- Taxonomy}{The four-class operational taxonomy (Normal/Jammer/Spoofer/Interference) and its mapping from RadioML 2018.01A's 24 classes is not clearly described.}
Please see our response to Comment~1.2. The mapping is listed in Table~I, together with the
preprocessing used to build the balanced 30,000-sample test set (Sec.~II).

\rc{Comment 4.2 --- Accuracy}{83.23\% accuracy is low for the stated mission. For a security-critical task (countering jammers/spoofers in microseconds), a $\sim$17\% error rate is substantial.}
We agree. The per-class results (Table~III) show where the errors come from:
\begin{itemize}
  \item Normal and Spoofer are recognized with F1 scores of 99\% and 95\%;
  \item nearly all errors are confusions between Jammer and Interference, because higher-order QAM
  modulations are split between these two groups and cannot be reliably distinguished from 128
  samples.
\end{itemize}
The limitation therefore lies in the class definition rather than in the hardware: the
accelerator reproduces the model bit-exactly, at 83.18\% compared with 83.26\% in FP32. This is
why we are moving to the injected-interference dataset described in our response to
Comment~1.2. In that dataset the classes correspond to physically distinct interference types.

\rc{Comment 4.3 --- Baseline fairness}{The GPU baseline comparison is methodologically unfair --- comparing FPGA datapath latency against batch-1 GPU end-to-end latency inflates the claimed advantage; fairer baselines (on-chip ARM cores, MCU-class processors, or comparable FPGA accelerators) should be used.}
We agree, and the GPU comparison has been removed. The revised paper uses the reviewer's
suggested baseline. The same integer model, implemented in C and compiled with \texttt{-O3},
runs on the ARM Cortex-A9 of the same Zynq device, and both implementations are measured on the
board with the same timer. The ARM takes 9.79~ms per inference; the accelerator takes
161.9~$\mu$s end to end, including data transfer. Both produce logits bit-identical to the golden
model (Table~IX). Comparable FPGA accelerators are contrasted in Table~V.

\rc{Comment 4.4 --- Clock frequency}{The 72.2~MHz clock frequency is unusually low for the Zynq-7000 platform, suggesting potential timing closure issues or a serialized datapath with limited parallelism.}
The reviewer's diagnosis was correct: the original datapath was both serialized and not
pipelined. Please see our response to Comment~3.2. The redesign performs $W^2{=}64$ MACs per
cycle and runs at 142.86~MHz in the complete system (175.4~MHz out of context).

\end{document}
